\begin{textsample}{POS Dim 5 – global\_north\_2024\_02 – Score 16.00 – global\_north\_2024\_02/105945058.txt}  \label{ex:f5_pos_238}
Israel has begun bombing Rafah, as part of a planned full-scale ground offensive on the city.
About 1.4 million of the 2.3 million people who live in Gaza are currently sheltering in Rafah in dire conditions, but with nowhere else safe to go.
Since Israel began its war on Gaza in response to the Hamas attacks on Israel on 7 October, Israel has killed more than 28,500 people, mostly women and children ; much of the civilian infrastructure has been destroyed and more than 80 percent of the population has been displaced.
Many of the displaced civilians have sought refuge in Rafah.
Despite those stark facts, it is against that background that the head of UN aid Martin Griffiths is warning that Israel ' s planned offensive in Rafah could lead to"slaughter", and that this"long-dreaded scenario is unraveling at alarming speed".
New Zealand ' s Prime Minister and Foreign Minister have"urged"Israel not to begin its ground offensive in Rafah, as part of"an overwhelming consensus of the \textbf{international} community".
Rt Hon Winston @ @ @ @ @ @ @ @ @ @ by indications that Israel is planning a ground offensive into Rafah"and"the humanitarian consequences of an offensive in Rafah would be appalling."We welcome these statements of concern.
They are a necessary indication of our government ' s attitude to the next phase of Israel ' s war on Gaza.
But they are not sufficient.
They will not reassure Palestinian refugee families sheltering in tents at the border with Egypt tonight.
They will not alter Israel ' s \textbf{genocidal} intention to exterminate and displace Palestinian civilians, under cover of a war on Hamas.
Advertisement - \textbf{scroll} to continue reading Are you getting our free newsletter?
The South \textbf{African} government has lodged an urgent request with the \textbf{International} \textbf{Court} of Justice to consider whether Israel ' s operations targeting Rafah are a breach of the \textbf{provisional} orders the \textbf{court} made in the \textbf{case} alleging \textbf{genocide} by Israel, and to order additional \textbf{provisional} orders to halt the mass killing in Rafah.
The Foreign Minister has noted that NZ regards the ICJ ' s decisions, including the @ @ @ @ @ @ @ @ @ @ the \textbf{Genocide} \textbf{Convention}, as binding.
The New Zealand Government should support South Africa ' s urgent request to the ICJ and take other concrete steps to \textbf{sanction} Israel for its failure to comply with \textbf{international} law.
If we truly want to hold Israel to account, the Israeli ambassador should be left in no doubt that, if the Rafah ground offensive goes ahead, diplomatic relations with Israel will cease."We are at the precipice of witnessing the mass slaughter of civilians in Rafah and the most concerted effort yet to depopulate Gaza.
These are flagrant violations of \textbf{international} law and the ICJ ' s orders.
New Zealand ' s response needs to \textbf{measure} up to the enormity of the situation at hand"Marilyn Garson, Alternative Jewish Voices.
Did you know \textbf{Scoop} has an Ethical Paywall?
If you ' re using \textbf{Scoop} for work, your organisation needs to pay a small \textbf{license} \textbf{fee} with \textbf{Scoop} \textbf{Pro}.
We think that ' s \textbf{fair}, because your organisation is benefiting from using our news resources.
In return, @ @ @ @ @ @ @ @ @ @ tools and keep \textbf{Scoop} free for personal use, because public access to news is important!

% matched lemmas: african, case, convention, court, fair, fee, genocidal, genocide, international, license, measure, pro, provisional, sanction, scoop, scroll
\end{textsample}
